\documentclass{article}

\usepackage{comment}

\def\foo#1{}

\begin{document}

Hello, world.

% Here's a normal comment. This will be removed.

\iffalse
  Sometimes \TeX\ authors will abuse \texttt{iffalse} as a
    commenting mechanism. This is accounted for and will be
    removed. 
\fi

\begin{comment}
  There are also explicit comment environments defined by,
  e.g., the `comment` package. We remove these, too.
\end{comment}

\foo{
  Sometimes \TeX\ authors will also create special macros for
  the purposes of commenting things out. We again remove these.
}

\begin{verbatim}
Other times, we don't want to remove comments, such as inside
of verbatim environments. So we protect against these cases.
% This comment won't be removed.
\end{verbatim} 

You might wonder why comment removal is so important to
me. The reason is that, in \TeX\ communities, comments are
rarely used as they are in programming or development
communities. Namely, \TeX\ users typically use comments to
\emph{remove} unwanted content rather than to comment on
wanted content.

For example, suppose you are told to correct the typo
'hommorphism' in the following text. Can you easily do so
with the comments?

% Let $f:X \to Y$ be a hommorphism.
\iffalse
  Let $f:X \to Y$ be a hommorphism.
\begin{comment}
  Let $f:X \to Y$ be a hommorphism.
\end{comment}
  Let $f:X \to Y$ be a hommorphism.
% \fi
  Let $f:X \to Y$ be a hommorphism.\fi
\foo{Let $f:X \to Y be %
  %
  %
  a hommorphism.\iffalse
%} Let $f:X \to Y$ be a hommorphism.
Let $f:X \to Y$ be a hommorphism.}
Let $f:X \to Y$ be a hommorphism.
%Let $f:X \to Y$ be a hommorphism.
% \fi



\end{document}
